
\documentclass[12]{article}
\usepackage{graphicx}
\usepackage[T1]{fontenc}
\usepackage{tgbonum}
\usepackage{amsmath}
\usepackage{amssymb}
\usepackage{booktabs}
\usepackage{float}
\usepackage{graphicx} % Required for inserting images
\usepackage{ulem}
\usepackage{hyperref}
\usepackage{fontawesome5}
\usepackage{dirtytalk}
\usepackage{xcolor}
\usepackage[affil-it]{authblk}

\usepackage{setspace}
\doublespacing

\usepackage{geometry}
\geometry{
 a4paper,
 total={170mm,257mm},
 left=20mm,
 top=20mm,
 }






\title{EK369E: Financial Markets\\ Term Structure and Interest Rate Risk\\ Problem Set (with solutions)}
\author{Nord University Business School}
\date{}

\begin{document}

\maketitle



\section*{Problem 1}
The table below summarises some information from the fixed-income market:
% Please add the following required packages to your document preamble:
% \usepackage{booktabs}
% \usepackage{graphicx}
\begin{table}[H]
\centering
\resizebox{0.3\textwidth}{!}{%
\begin{tabular}{@{}llll@{}}
\toprule
Bond & \begin{tabular}[c]{@{}l@{}}Maturity\\ (years)\end{tabular} & \begin{tabular}[c]{@{}l@{}}Coupon rate\\ (\%)\end{tabular} & \begin{tabular}[c]{@{}l@{}}Yield\\ (\%)\end{tabular} \\ \midrule
A    & 5.5                                                        & 0                                                          & 2.5                                                  \\
B    & 2.5                                                        & 4.0                                                        & 2.0                                                  \\ \bottomrule
\end{tabular}%
}
\label{tab:my-table}
\end{table}

\begin{enumerate}
    \item[a)] Calculate the duration and the modified duration for the two bonds.
\end{enumerate}


\color{blue}
\subsection*{Solution 1a}
The duration of Bond A is equal to its time to maturity i.e. 5.5 yrs (because it's a zero-coupon bond)\\
The modified duration of Bond A: $D^* = ${\Large $\frac{5.5}{1.025}$} = \uuline{5.37}\\

\noindent\textbf{Working for Bond B:} Assume annual coupon payments and a face value of NOK 100. The annual coupon payment is $\text{Coupon rate}\times\text{Face value}=0.04\times100=\text{NOK }4$. With 2.5 years remaining, the bond pays NOK 4 in 0.5 years, NOK 4 in 1.5 years, and NOK 104 in 2.5 years. The final payment includes the face value and the last coupon. \par\noindent In Macaulay's duration formula, the term $CF_t(1+r)^{-t}$ gives the present value of each payment: \[PV_t=CF_t(1+r)^{-t}=\frac{CF_t}{(1+r)^t}.\] Here, $CF_t$ is the cash payment received at time $t$, $r$ is the bond's annual yield, $t$ is the time until payment in years, and $PV_t$ is the value of that payment today. For Bond B, the given yield is $r=0.02$. For example, the first payment has present value \[PV_{0.5}=\frac{4}{(1+0.02)^{0.5}}\approx3.9606.\] The table repeats this calculation for each payment. The total of the PV column gives the bond price, while the total of the $PV_t\times t$ column gives the numerator used in the duration calculation below. \par\noindent
\\

\begin{table}[H]
    \centering
    {\color{blue}\begin{tabular}{ccccc}
    \toprule
        Time $t$ & Cash flow $CF_t$ & $(1+r)^t$, with $r=2\%$ & $PV_t$ & $PV_t\times t$\\\midrule
        0.5 & 4 & 1.0100 & 3.9606 & 1.9803\\
        1.5 & 4 & 1.0301 & 3.8829 & 5.8244\\
        2.5 & 104 & 1.0508 & 98.9767 & 247.4417\\\midrule
            &     &        & 106.8202 &255.2464\\ \bottomrule
    \end{tabular}}
    \label{tab:my_label}
\end{table}

\noindent
\\
Duration of Bond B: {\Large $\frac{255.2464}{106.8202}$} = \uuline{2.39}\\
\\
Modified Duration of Bond B: {\Large $\frac{2.39}{1.02}$} = \uuline{2.34}\\






\color{black}
\begin{enumerate}
    \item[b)] Which of the two bonds will have the smallest price change for a given change in the interest rate? Give a short explanation.
\end{enumerate}

\color{blue}
\subsection*{Solution 1b}
For any 1\% increase in interest rates:\\
Bond A will experience a 5.4\% reduction in its price:$-5.37 \times 0.01 =  -0.054$\\
Bond B will experience a 2.3\% reduction in its price:$-2.34 \times 0.01 =  -0.023$\\
Bond B has a shorter effective maturity (duration) and a higher coupon rate, therefore it is less vulnerable to future fluctuation in interest rates.







\color{black}
\begin{enumerate}
    \item[c)] If a fund manager invests (for face value) NOK 10 million in bond A and (for face value) NOK 16 million in bond B, what will the duration of the portfolio be?
\end{enumerate}

\color{blue}
\subsection*{Solution 1c}
Price of Bond A: {\Large $\frac{100}{(1.025)^{5.5}}$} = 87.3\\
Market Value of Bond A: NOK 10 million {\Large$ \times \frac{87.3}{100} $} =  NOK 8.73 million\\
\\
Market Value of Bond B: NOK 16 million {\Large$ \times \frac{106.82}{100}$} = NOK 17.09 million\\
\\
Total Value of the Bond Portfolio: NOK 25.82 million\\
\\
Weight of Bond A in the portfolio: {\Large $\frac{8.73 million }{25.82 million}$} = 0.34\\
\\
Weight of Bond B in the portfolio: {\Large $\frac{17.09 million }{25.82 million}$} = 0.66\\
\\
Duration of the portfolio: $(0.34 \times 5.5) + (0.66 \times 2.39) = $ \uuline{3.44}



{\Large}



\noindent
\\
\color{black}
\color{black}
\begin{enumerate}
    \item[d)] Calculate the duration of the portfolio one year later, given that the yield is unchanged for both bonds.
\end{enumerate}

\color{blue}
\subsection*{Solution 1d}
One year later\\
The duration of Bond A = 4.5 years\\
Duration of Bond B: {\Large $\frac{153.4146}{104.9168}$} = 1.46 years\\

% Please add the following required packages to your document preamble:
% \usepackage{booktabs}
% \usepackage{graphicx}
\begin{table}[H]
\centering
\resizebox{0.4\textwidth}{!}{%
\color{blue}\begin{tabular}{@{}lllll@{}}
\toprule
t   & CF  & $(1+r)^t$, $r=2\%$     & PV       & $PV*t$     \\ \midrule
0.5 & 4   & 1.0100 & 3.96059  & 1.980295 \\
1.5 & 104 & 1.0301 & 100.9562 & 151.4343 \\ \midrule
    &     &        & 104.9168 & 153.4146 \\ \bottomrule
\end{tabular}%
}

\label{tab:my-table}
\end{table}

\noindent\textbf{Duration of the portfolio one year later:} Assume the coupon received during the year is withdrawn from the portfolio and the face-value holdings remain unchanged. One year later, Bond A has 4.5 years remaining and Bond B has 1.5 years remaining. We first calculate their updated prices using the bond-pricing formula:
\[P=\sum_t\frac{CF_t}{(1+r)^t},\]
where $P$ is the price per NOK 100 face value, $CF_t$ is each remaining payment per NOK 100 face value, $r$ is the bond's annual yield, and $t$ is the time until payment measured from the new valuation date. The yields remain 2.5\% for Bond A and 2\% for Bond B. Therefore,
\[P_1^A=\frac{100}{(1+0.025)^{4.5}}\approx89.4834,\]
\[P_1^B=\frac{4}{(1+0.02)^{0.5}}+\frac{104}{(1+0.02)^{1.5}}\approx104.9168.\]
The price of Bond B also equals the total of the PV column in the table above.
\par\noindent The market value of each holding is calculated as:
\[\text{Market value}=\text{Face value held}\times\frac{\text{Price per NOK 100 face value}}{100}.\]
Using the face-value holdings from part (c), the market values below are expressed in NOK million: \[MV_A=10\times\frac{89.4834}{100}\approx8.9483,\] \[MV_B=16\times\frac{104.9168}{100}\approx16.7867.\]
Portfolio duration is the market-value-weighted average of the individual bond durations:
\[D_P=\frac{MV_A D_A+MV_B D_B}{MV_A+MV_B},\]
where $MV_A$ and $MV_B$ are the market values and $D_A$ and $D_B$ are the individual durations at the same date. Using $D_A=4.5$ years and $D_B\approx1.4623$ years from the calculations above, before rounding Bond B's duration to 1.46:
\[D_{P,1}=\frac{8.9483\times4.5+16.7867\times1.4623}{8.9483+16.7867}\approx\uuline{2.52\text{ years}}\] \par\noindent


\noindent
\\
\color{black}
\color{black}
\begin{enumerate}
    \item[e)] What is duration drift, and how big is the duration drift for the two bonds?
\end{enumerate}

\color{blue}
\subsection*{Solution 1e}
\noindent With unchanged yields, duration need not decline by exactly one year when one year passes. Here, duration drift measures the difference between the time elapsed and the reduction in duration:
\[
\text{Duration drift}=\text{Time elapsed}-(\text{Initial duration}-\text{New duration}).
\]
\noindent\textbf{Bond A:}
Its duration falls by exactly one year, matching the passage of time:
\[
\text{Duration drift}=1-(5.5-4.5)=1-1=\text{\uuline{0 years}}.
\]
\noindent\textbf{Bond B:}
Its duration falls by approximately 0.93 years, which is less than the one year that has passed:
\[
\text{Duration drift}=1-(2.39-1.46)=1-0.93\approx\text{\uuline{0.07 years}}.
\]
\noindent Bond B's duration therefore declines by approximately 0.07 years less than its time to maturity.\\

% Please add the following required packages to your document preamble:
% \usepackage{booktabs}
% \usepackage{graphicx}
\begin{table}[H]
\centering
\resizebox{0.9\textwidth}{!}{%
\color{blue}\begin{tabular}{@{}lllllll@{}}
\toprule
Bond & Initial duration & New duration & $\Delta$Duration  & Initial Maturity & New Maturity & $\Delta$ Maturity  \\ \midrule
A    & 5.5              & 4.5          & 1                  & 5.5              & 4.5          & 1                  \\
B    & 2.39             & 1.46         & 0.93               & 2.5              & 1.5          & 1                  \\ \bottomrule
\end{tabular}%
}
\label{tab:my-table}
\end{table}







\color{black}
\section*{Problem 2}
In the fixed-income market, four different zero-coupon bonds are traded, all with a face value of NOK 100. The table below summarizes the relevant data:

% Please add the following required packages to your document preamble:
% \usepackage{booktabs}
% \usepackage{graphicx}
\begin{table}[H]
\centering
\resizebox{0.25\textwidth}{!}{%
\begin{tabular}{@{}lll@{}}
\toprule
Bond & Price & \begin{tabular}[c]{@{}l@{}}Years to\\  Maturity\end{tabular} \\ \midrule
A    & 97.75 & 1                                                            \\
B    & 95.00 & 2                                                            \\
C    & 92.00 & 3                                                            \\
D    & 89.00 & 4                                                            \\ \bottomrule
\end{tabular}%
}
\label{tab:my-table}
\end{table}

\begin{enumerate}
    \item[a)] Calculate the rates on the spot yield curve i.e.  $s_1, s_2, s_3$ and $s_4$.
\end{enumerate}

\color{blue}
\subsection*{Solution 2a}
The 1-yr spot rate: 97.75 = {\Large $ \frac{100}{1 + r_1}$} \\
$r_1 =$ {\Large $ \frac{100}{97.75}$} - 1 = \uuline{2.30\%}\\
\\
The 2-yr spot rate: 95.00 = {\Large $ \frac{100}{(1 + r_2)^2}$} \\
$r_2 =$ {\Large $ \sqrt[2]{\frac{100}{95.0}}$} - 1 = \uuline{2.60\%}\\
\\
The 3-yr spot rate: 92.00 = {\Large $ \frac{100}{(1 + r_3)^3}$} \\
$r_3 =$ {\Large $ \sqrt[3]{\frac{100}{92.0}}$} - 1 = \uuline{2.82\%}\\
\\
The 4-yr spot rate: 89.00 = {\Large $ \frac{100}{(1 + r_4)^4}$} \\
$r_4 =$ {\Large $ \sqrt[4]{\frac{100}{89.0}}$} - 1 = \uuline{2.96\%}\\




\color{black}
\begin{enumerate}
    \item[b)] A three-year coupon bond has a face value of NOK 100 and an annual coupon rate of 3\%.
    \begin{enumerate}
        \item[i] Based on the yield structure above, what is the market price of the coupon bond?
        \item[ii] Find the yield to maturity for the coupon bond.
        \item[iii] Find the duration of the coupon bond.
    \end{enumerate}
\end{enumerate}


\color{blue}
\subsection*{Solution 2b}
i.\\
Current market price of the 3\% coupon bond with three years to maturity:\\
$P_0 =$ {\Large $\frac{3}{1.0230} + \frac{3}{1.0260^2} + \frac{103}{1.0282^3}$} = \uuline{100.54}\\
\\
ii.\\
100.54 = {\Large $\frac{3}{1+YTM} + \frac{3}{(1+YTM)^2} + \frac{103}{(1+YTM)^3}$}\\
YTM $\approx$ \uuline{2.81\%}\\
\\
iii.\\
Duration of the bond: {\Large $\frac{292.944}{100.539}$} = \uuline{2.91 yrs}\\

\begin{table}[H]
    \centering
    {\color{blue}\begin{tabular}{ccccc}
    \toprule
        time & Cashflow & $(1+\mathrm{YTM})^t$ & PV & $PV \times t$\\\midrule
        1 & 3 & 1.0281 & 2.9180 & 2.9180\\
        2 & 3 & 1.0570 & 2.8382 & 5.6765\\
        3 & 103 & 1.0867 & 94.7832 & 284.3495\\\midrule
         & &        & 100.5394 & 292.944\\ \bottomrule
    \end{tabular}}
    \label{tab:my_label}
\end{table}




\color{black}
\begin{enumerate}
    \item[c)] Calculate the implied one-year forward rates in year 2, year 3 and year 4.
\end{enumerate}

\color{blue}
\subsection*{Solution 2c}
%The one-year forward rate in year 2: \\
$f_{1,2} = ${\Large $\frac{1.0260^2}{1.0230}$} - 1 = \uuline{2.89\%}\\
\\
$f_{2,3} = ${\Large $\frac{1.0282^3}{1.0260^2}$} - 1 = \uuline{3.26\%}\\
\\
$f_{3,4} = ${\Large $\frac{1.0296^4}{1.0282^3}$} - 1 = \uuline{3.37\%}\\


\color{black}
\begin{enumerate}
    \item[d)] Assuming that the spot yield curve evolves in line with the pure expectation hypothesis, compute the one-year holding period return on this bond.
\end{enumerate}

\color{blue}
\subsection*{Solution 2d}
The yield curve one-year later:\\
\\
1-yr spot rate: $ r_1$ = $ f_{1,2}$ = 2.89\%\\
\\
2-yr spot rate: $(1 + r_2)^2$ = $(1+r_1) (1 + f_{2,3})$\\
$r_2 =$ ($\sqrt[2]{1.0289 \times 1.0326}$) - 1 = \uuline{3.08\%}\\
\\
3-yr spot rate: $(1 + r_3)^3$ = $(1+r_1) (1 + f_{2,3}) (1 + f_{3,4})$\\
$r_3 =$ ($\sqrt[3]{1.0289 \times 1.0326 \times 1.0337}$) - 1 = \uuline{3.18\%}\\
\\
The price of the 3\% coupon bond one-year later:\\
$P_1 =$ {\Large $\frac{3}{1.0289} + \frac{103}{1.0308^2}$} = \uuline{99.86}\\
\\
Holding period return: HPR = {\Large $\frac{99.86 - 100.54 + 3}{100.54}$} = \uuline{2.30\%}



\color{black}
\section*{Problem 3}
A government bond, Bond 1, with a face value of 100 and coupon rate of 4.50\% will mature in four years. An investor is looking to invest NOK 5 million for a period of two years. The investor considers various strategies in the interest rate market. In addition to Bond 1, the investor may invest in government Bond 2. Bond 2’s coupon rate is also 4.50\%, but the time to maturity is two years. In the financial press, you see that the term structure (yield curve) for government bonds with a coupon rate of 4.50\% is:
% Please add the following required packages to your document preamble:
% \usepackage{graphicx}
\begin{table}[H]
\centering
\resizebox{0.4\textwidth}{!}{%
\begin{tabular}{lllll}
Year       & 1    & 2    & 3    & 4    \\\midrule
Yield (\%) & 4.00 & 4.20 & 4.50 & 4.90
\end{tabular}%
}
\label{tab:my-table}
\end{table}
\begin{enumerate}
    \item[a] Assume that the yield curve is stable, calculate the geometric return for:
    \begin{enumerate}
        \item[i)] Buy bond 1, and sell it after two years
        \item[ii)] Buy bond 2, and hold it to maturity.
    \end{enumerate}
    
\end{enumerate}

\color{blue}
\subsection*{Solution 3a}
\textbf{i) Buy Bond 1 and sell it after two years}\\
\\
Price today @ t = 0:\\
$P_0^{1} = ${\Large $\frac{4.5}{1.049} + \frac{4.5}{1.049^2} + \frac{4.5}{1.049^3} + \frac{104.5}{1.049^4}$} = 98.58\\
\\
Price one year later @ t = 1:\\
$P_1^{1} = ${\Large $\frac{4.5}{1.045} + \frac{4.5}{1.045^2} + \frac{104.5}{1.045^3}$} = 100\\
\\
Price two years later @ t = 2:\\
$P_2^{1} = ${\Large $\frac{4.5}{1.042} + \frac{104.5}{1.042^2}$} = 100.56\\
\\
Holding Period Return: \\
$HPR_1^1 =$ {\Large $\frac{100-98.58 + 4.5}{98.58}$} = 6\%\\
\\
$HPR_2^1 =$ {\Large $\frac{100.56 - 100 + 4.5}{100}$} = 5.06\%\\
\\
Geometric average return over investment period: $r_G^1 =$ $\sqrt{1.06 \times 1.0506} - 1$ = \uuline{5.53\%}
\\

\noindent
\\
\textbf{ii) Buy bond 2, and hold it to maturity}\\
%The geometric average return on Bond 2 which is held to maturity (assuming that the yield curve remains stable) is equal to the bond's yield to maturity i.e 4.20\%\\
Price today @ t = 0:\\
$P_0^{2} = ${\Large $\frac{4.5}{1.042} + \frac{104.5}{1.042^2}$} = 100.56\\
\\
Price one year later @ t = 1:\\
$P_1^{2} = ${\Large $\frac{104.5}{1.040}$} = 100.48\\
Holding Period Return: $HPR_1 =$ {\Large $\frac{100.48 - 100.56 + 4.5}{100.56}$} = 4.39\%\\
\\
Price two years later @ t = 2:\\
$P_2^{2} = $ 100 (bond is redeemed at par value)\\
Holding Period Return: 
$HPR_2^2 =$ {\Large $\frac{100 - 100.48 + 4.5}{100.48}$} = 4.00\%\\
\\
Geometric average return over investment period: $r_G^2 =$ $\sqrt{1.0439 \times 1.04} - 1$ = \uuline{4.2\%}

\noindent
\\
\color{black}
\begin{enumerate}
    \item[b)] Explain why there are differences between the two strategies in a)   
\end{enumerate}


\color{blue}
\subsection*{Solution 3b}
Investment strategy 1 is called rolling down or riding the yield curve while Investment strategy 2 is known as Maturity Matching. When the yield curve is upward-sloping and stable, investment strategy 1 offers superior returns over strategy 2.





\color{black}
\section*{Problem 4}
In the fixed-income market, we have the following information for zero-coupon bonds:

\begin{table}[H]
    \centering
    \begin{tabular}{cccc}
        Maturity (years) & 1 & 2 & 3\\\midrule
        Yield to maturity (\%) & 4.3 & 4.0 & 3.8\\
    \end{tabular}
    \label{tab:my_label}
\end{table}

\noindent
You are analysing three bonds, all maturing in three years, but with different coupon rates. Bond A has a coupon rate of 8\%, bond B has a coupon rate of 4\%, while bond C is a zero coupon bond.

\begin{enumerate}
    \item [a)] Calculate the price of the three bonds.
\end{enumerate}

\color{blue}
\subsection*{Solution 4a}
Price of Bond A:\\
$P_0^A =$ {\Large $\frac{8}{1.043} + \frac{8}{1.04^2} + \frac{108}{1.038^3}$} = \uuline{111.63}\\
\\
Price of Bond B:\\
$P_0^B =$ {\Large $\frac{4}{1.043} + \frac{4}{1.04^2} + \frac{104}{1.038^3}$} = \uuline{100.52}\\
\\
Price of Bond C:\\
$P_0^C =$ {\Large $\frac{100}{1.038^3}$} = \uuline{89.41}\\
\\




\color{black}
\begin{enumerate}
    \item [b)] Without making calculations, briefly explain which bond has the lowest yield to maturity and which bond has the lowest duration.
\end{enumerate}

\color{blue}
\subsection*{Solution 4b}
\noindent Bond C has the lowest yield to maturity, 3.8\%, because its only payment occurs at its three-year maturity. Its yield therefore equals the three-year zero-coupon yield. Bonds A and B also make earlier coupon payments, which are valued using the higher one-year and two-year yields on the given downward-sloping yield curve. Their yields to maturity are therefore higher than Bond C's, with Bond A having the highest yield to maturity.
\\
Bond A has the lowest duration. With the highest coupon rate, Bond A offers more of its total cash flow at earlier periods in the bond's life. In contrast, Bond C has the highest duration since all of its cash flow is offered at maturity.



\noindent
\\
\color{black}
You have an investment horizon of two years and are considering investing in either bond A or bond C. You are considering two different strategies: \textbf{i)} Buy bond A today, reinvest the coupon payment in one year in the same security and sell in two years. \textbf{ii)} Buy bond C today and sell it in two years. You assume that the yield curve given in the table above is stable.
\begin{enumerate}
    \item [c)] Calculate the annual geometric average return for each of the two investment strategies.
\end{enumerate}

\color{blue}
\subsection*{Solution 4c}
i) Buy bond A today, reinvest the coupon payment in one year in the same security, and sell in two years\\
\\
$P_0^A =$ \uuline{111.63}\\
$P_1^A =$ {\Large $\frac{8}{1.043} + \frac{108}{1.040^2}$} = \uuline{107.52}\\
$P_2^A =$ {\Large $\frac{108}{1.043}$} = \uuline{103.55}\\

\noindent\textbf{Reinvestment of the first coupon:}
The first coupon is reinvested in the same bond at the end of year 1. It therefore earns Bond A's holding period return from year 1 to year 2:
\[
HPR_2=\frac{P_2^A-P_1^A+C}{P_1^A},
\]
where $P_1^A$ is the bond price at the end of year 1, $P_2^A$ is its price at the end of year 2, and $C$ is the annual coupon payment.
Using $P_1^A=107.52$, $P_2^A=103.55$, and $C=8$:
\[
HPR_2=\frac{103.55-107.52+8}{107.52}\approx0.0374814.
\]
The value of the reinvested first coupon at the end of year 2 is:
\[
\text{Reinvested coupon value}=C(1+HPR_2).
\]
Therefore:
\[
\text{Reinvested coupon value}=8(1+0.0374814)\approx8.30.
\]

\begin{table}[H]
    \centering
    \resizebox{0.6\textwidth}{!}{%
    {\color{blue}\begin{tabular}{cll}
    \toprule
         & Description & Cash flow \\\midrule
         & 1st coupon received at $t=1$ and reinvested & $8(1+0.0374814)\approx8.30$\\
         & 2nd coupon received at t = 2  & 8 \\
         & Proceeds from bond sale  & 103.55 \\\bottomrule
         & Settlement amount & 119.85\\\bottomrule
    \end{tabular}}
    }
    \label{tab:my_label}
\end{table}

\noindent The settlement amount consists of the reinvested first coupon, the second coupon, and the proceeds from selling the original bond:
\[
V_2=C(1+HPR_2)+C+P_2^A.
\]
Using the amounts in the table:
\[
V_2=8.30+8+103.55=119.85.
\]
The annual geometric average return is calculated using:
\[
r_G=\left(\frac{V_T}{V_0}\right)^{1/T}-1,
\]
where $V_0$ is the initial investment, $V_T$ is the settlement amount after $T$ years, and $r_G$ is the annual geometric average return.
Here, $V_0=111.63$, $V_T=V_2=119.85$, and $T=2$:
\[
r_G=\sqrt{\frac{119.85}{111.63}}-1\approx\text{\uuline{3.62\%}}.
\]

\noindent
\\
\textbf{ii) Buy bond C today and sell it in two years}\\
\\
Price of Bond C:\\
$P_0^C =$ {\Large $\frac{100}{1.038^3}$} = \uuline{89.41}\\
\\
$P_1^C =$ {\Large $\frac{100}{1.040^2}$} = \uuline{92.46}\\
\\
$P_2^C =$ {\Large $\frac{100}{1.043}$} = \uuline{95.88}\\
\\
Holding Period Return (HPR):
$HPR_1 =$ {\Large $\frac{92.46 - 89.41}{89.41}$} = 3.4\%\\
\\
$HPR_2 =$ {\Large $\frac{95.88 - 92.46}{92.46}$} = 3.7\%\\
\\
Geometric return: $\sqrt{1.034 \times 1.037} - 1$ = \uuline{3.55\%}


\noindent
\\
\color{black}
\begin{enumerate}
    \item [d)] Which alternative is best? Briefly explain why.
\end{enumerate}

\color{blue}
\subsection*{Solution 4d}
\noindent\textbf{Strategy 1 (Bond A) is better.}
The results in part (c) show an annual geometric average return of approximately 3.62\% for Bond A, compared with 3.55\% for Bond C.

\par\noindent On the unchanged, downward-sloping yield curve, shorter maturities have higher yields. As the bonds' remaining maturities shorten, their yields therefore rise. Bond A has a lower duration than Bond C, making its price less sensitive to this effect.

\par\noindent Neither strategy is maturity matching, because both involve buying a three-year bond and selling it after two years, before maturity.


\noindent

\color{black}
\begin{enumerate}
    \item [e)] Is there an alternative investment strategy in this market that would be clearly better than the two mentioned here?
\end{enumerate}

\color{blue}
\subsection*{Solution 4e}
\noindent\textbf{A better alternative is a roll-over strategy.}
Buy a one-year zero-coupon bond today (Bond D). When it matures after one year, reinvest all its redemption proceeds in another one-year zero-coupon bond (Bond E), which matures at the end of year 2.

\par\noindent Since the yield curve remains unchanged, the one-year yield is 4.3\% at both purchase dates. The 4.0\% yield applies to a bond with two years remaining, not to a one-year bond purchased in the second year.

\par\noindent The price of a zero-coupon bond is:
\[
P=\frac{FV}{(1+r)^T},
\]
where $P$ is the bond price, $FV$ is its face value, $r$ is its annual yield, and $T$ is its remaining maturity in years.
For each purchase, $FV=100$, $r=0.043$, and $T=1$:
\[
P_0^D=P_1^E=\frac{100}{(1+0.043)^1}\approx95.88.
\]
Each bond is redeemed for NOK 100 after its one-year holding period.

\par\noindent The holding period return is:
\[
HPR=\frac{P_{\mathrm{end}}-P_{\mathrm{start}}+C}{P_{\mathrm{start}}},
\]
where $P_{\mathrm{start}}$ is the purchase price, $P_{\mathrm{end}}$ is the sale or redemption value, and $C$ is the coupon received during the holding period.
Both bonds are zero-coupon bonds, so $C=0$:
\[
HPR_1=HPR_2=\frac{100-95.88+0}{95.88}\approx0.043=4.30\%.
\]

\par\noindent The annual geometric average return over the two years is:
\[
r_G=\sqrt{(1+HPR_1)(1+HPR_2)}-1,
\]
where $HPR_1$ and $HPR_2$ are the returns in the first and second years.
Therefore:
\[
r_G=\sqrt{(1+0.043)(1+0.043)}-1=\text{\uuline{4.30\%}}.
\]
\noindent This exceeds the approximately 3.62\% return from Bond A and 3.55\% return from Bond C calculated in part (c).




\noindent
\\
\color{black}
You put together a portfolio consisting of a face value of NOK 100 million in bond A and a face value of NOK 91.9 million in a one-year zero-coupon bond D.
\begin{enumerate}
    \item [f)] What is the duration of this portfolio?
\end{enumerate}

\color{blue}
\subsection*{Solution 4f}
\noindent\textbf{Bond D:}
For a zero-coupon bond, duration equals its remaining maturity:
\[
D_D=T_D=1\text{ year}.
\]

\par\noindent\textbf{Bond A:}
To calculate Macaulay duration, first find the bond's yield to maturity using:
\[
P_0=\sum_{t=1}^{T}\frac{CF_t}{(1+r)^t},
\]
where $P_0$ is the current bond price, $CF_t$ is the payment received at time $t$, $r$ is the annual yield to maturity, and $T$ is the remaining maturity in years.

\par\noindent Using Bond A's price from part (a), keeping extra decimal places:
\[
111.634289=\frac{8}{1+r}+\frac{8}{(1+r)^2}+\frac{108}{(1+r)^3}.
\]
Solving for the yield using a financial calculator gives:
\[
r_A\approx0.0382177449=3.82177449\%.
\]
Use this same yield for all payments in the Macaulay duration calculation.

\par\noindent Macaulay's duration formula is:
\[
D=\sum_{t=1}^{T}t\left[\frac{CF_t(1+r)^{-t}}{P_0}\right],
\]
where $D$ is duration in years.
For the table below, calculate each payment's present value using:
\[
PV_t=\frac{CF_t}{(1+r_A)^t}.
\]
Then multiply each present value by its payment time $t$.

\begin{table}[H]
    \centering
    {\color{blue}
    \begin{tabular}{ccccc}
    \toprule
    Time $t$ & Cash flow $CF_t$ & $(1+r_A)^t$ & $PV_t$ & $PV_t\times t$\\
    \midrule
    1 & 8   & 1.0382 & 7.7055  & 7.7055\\
    2 & 8   & 1.0779 & 7.4219  & 14.8437\\
    3 & 108 & 1.1191 & 96.5069 & 289.5207\\
    \midrule
      &     &        & 111.6343 & 312.0700\\
    \bottomrule
    \end{tabular}}
\end{table}

\noindent Dividing the total of the last column by the bond price gives:
\[
D_A=\frac{\sum_t t\,PV_t}{P_0^A}
=\frac{312.0700}{111.6343}
\approx\text{\uuline{2.7955 years}}.
\]

\noindent
Purchase amount for Bond A: NOK 100 million $\times \frac{111.634}{100}$ = NOK 111.634 million\\
\\
Purchase amount for Bond D: NOK 91.9 million $\times \frac{95.88}{100}$ = NOK 88.114 million\\
\\
Value of Bond Portfolio: NOK 111.634 million + NOK 88.114 million = NOK 199.748 million\\
\\
\noindent The portfolio weights are calculated using market values:
\[
w_A=\frac{MV_A}{MV_A+MV_D},
\qquad
w_D=\frac{MV_D}{MV_A+MV_D},
\]
where $MV_A$ and $MV_D$ are the market values of the holdings in Bonds A and D, and $w_A$ and $w_D$ are their portfolio weights.

\par\noindent Using the market values calculated above, expressed in NOK million:
\[
w_A=\frac{111.634}{199.748}\approx0.5589,
\]
\[
w_D=\frac{88.114}{199.748}\approx0.4411.
\]

\par\noindent Portfolio duration is the market-value-weighted average of the individual durations:
\[
D_P=w_A D_A+w_D D_D,
\]
where $D_A$ and $D_D$ are the durations of Bonds A and D.
Therefore:
\[
D_P=(0.5589\times2.7955)+(0.4411\times1)
\approx\text{\uuline{2.00 years}}.
\]


\color{black}
\section*{Problem 5}
A newly issued bond has a maturity of 10 years and pays a 7\% coupon rate annually. The bond sells at par value.
\begin{enumerate}
    \item[a)] What is the duration of the bond?
\end{enumerate}

\color{blue}
\subsection*{Solution 5a}

\begin{table}[H]
    \centering
    {\color{blue}\begin{tabular}{ccccc}
    \toprule
        t & CF & $(1+r)^t$, $r=7\%$  & PV & $PV\cdot t$\\\midrule
        1 & 7 & 1.07 & 6.5421 & 6.5421 \\ 
        2 & 7 & 1.1449 & 6.1141 & 12.2281\\
        3 & 7 & 1.2250 & 5.7141 & 17.1423\\
        4 & 7 & 1.3108 & 5.3403 & 21.3611 \\
        5 & 7 & 1.4026 & 4.9909 & 24.9545\\
        6 & 7 & 1.5007 & 4.6644 & 27.9864\\
        7 & 7 & 1.6058 & 4.3592 & 30.5147\\
        8 & 7 & 1.7182 & 4.0741 & 32.5925\\
        9 & 7 & 1.8385 & 3.8075 & 34.2678\\
        10 & 107 & 1.9672 & 54.3934 & 543.9337 \\\midrule
           &     &         & 100 & 751.5232\\\bottomrule
    \end{tabular}}
    \label{tab:my_label}
\end{table}

\noindent
\\
Bond Duration: {\Large $\frac{751.5232}{100}$} = \uuline{7.52 yrs}


\noindent
\\
or 
% Please add the following required packages to your document preamble:
% \usepackage{booktabs}
% \usepackage{graphicx}
\begin{table}[H]
\centering
\resizebox{0.5\textwidth}{!}{%
\color{blue}\begin{tabular}{@{}llllll@{}}
\toprule
t  & CF  & $(1+r)^t$, $r=7\%$      & PV       & $w = \frac{PV_t}{P_0}$       & $w\cdot t$               \\ \midrule
1  & 7   & 1.07     & 6.542056 & 0.065421 & 0.065421          \\
2  & 7   & 1.1449   & 6.114071 & 0.061141 & 0.122281          \\
3  & 7   & 1.225043 & 5.714085 & 0.057141 & 0.171423          \\
4  & 7   & 1.310796 & 5.340266 & 0.053403 & 0.213611          \\
5  & 7   & 1.402552 & 4.990903 & 0.049909 & 0.249545          \\
6  & 7   & 1.50073  & 4.664396 & 0.046644 & 0.279864          \\
7  & 7   & 1.605781 & 4.359248 & 0.043592 & 0.305147          \\
8  & 7   & 1.718186 & 4.074064 & 0.040741 & 0.325925          \\
9  & 7   & 1.838459 & 3.807536 & 0.038075 & 0.342678          \\
10 & 107 & 1.967151 & 54.39337 & 0.543934 & 5.439337          \\ \midrule
   &     &          & 100      &          & \textbf{7.515232} \\ \bottomrule
\end{tabular}%
}

\label{tab:my-table}
\end{table}




\noindent
\\
\color{black}
\begin{enumerate}
    \item[b)] What is the convexity of the bond?
\end{enumerate}

\color{blue}
\subsection*{Solution 5b}

% Please add the following required packages to your document preamble:
% \usepackage{booktabs}
% \usepackage{graphicx}
\begin{table}[H]
\centering
\resizebox{0.5\textwidth}{!}{%
\color{blue}\begin{tabular}{@{}llllll@{}}
\toprule
t  & CF  & $(1+r)^t$, $r=7\%$     & PV      & $t^2 + t$ & $PV \cdot(t^2 + t)$ \\ \midrule
1  & 7   & 1.07   & 6.5421  & 2                        & 13.08411                      \\
2  & 7   & 1.1449 & 6.1141  & 6                        & 36.68443                      \\
3  & 7   & 1.2250 & 5.7141  & 12                       & 68.56902                      \\
4  & 7   & 1.3108 & 5.3403  & 20                       & 106.8053                      \\
5  & 7   & 1.4026 & 4.9909  & 30                       & 149.7271                      \\
6  & 7   & 1.5007 & 4.6644  & 42                       & 195.9046                      \\
7  & 7   & 1.6058 & 4.3592  & 56                       & 244.1179                      \\
8  & 7   & 1.7182 & 4.0741  & 72                       & 293.3326                      \\
9  & 7   & 1.8385 & 3.8075  & 90                       & 342.6783                      \\
10 & 107 & 1.9672 & 54.3934 & 110                      & 5983.271                      \\ \midrule
   &     &        &         &                          & 7434.175                      \\ \bottomrule
\end{tabular}%
}
\label{tab:my-table}
\end{table}

\noindent
Convexity of the bond: {\Large $\frac{1}{100} \cdot \frac{1}{1.07^2} $} $\cdot 7434.175$ = \uuline{64.93}





\noindent
\\
\color{black}
If the interest rate in the market is expected to increase from 7\% to 8\%;
\begin{enumerate}
    \item[c)] What new price is predicted by the modified duration rule 
\end{enumerate}

\color{blue}
\subsection*{Solution 5c}
Modified duration: {\Large $\frac{7.52}{1.07}$} = 7.02\\
price risk: $-7.02 \cdot 1\% = -7.02\%$ \\
New price predicted by the bond's duration: NOK 100$ \times 0.9298$ = \uuline{NOK 92.98}


\noindent
\\
\color{black}
If interest rates are expected to increase from 7\% to 8\% 
\begin{enumerate} 
    \item[d)] What new price is predicted by the convexity rule 
\end{enumerate}

\color{blue}
\subsection*{Solution 5d}
Price risk: $(-7.02 \times 0.01) + (0.5 \times 64.93 \times 0.01^2) = -0.067$ or -6.7\% \\
New price predicted by the duration and convexity of the bond: NOK 100$ \times 0.9330$ = \uuline{NOK 93.30}





\noindent
\\
\color{black}
\begin{enumerate}
    \item[e)] Find the actual price assuming that interest rates immediately increase from 7\% to 8\% (time to maturity is still  10 years) 
\end{enumerate}

\color{blue}
\subsection*{Solution 5e}
$P_0 = $ {\Large $7[\frac{1 - (1.08^{-10})}{0.08}] + \frac{100}{(1.08)^{10}}$} = \uuline{93.29}\\




\color{black}
\section*{Problem 6}
A 6\% coupon bond paying interest annually has a modified duration of 10-yrs, sells for NOK 80 and is priced at a yield to maturity of 8\%. If the YTM increases to 9\%, what is the predicted change in price based on the bond's duration

\color{blue}
\subsection*{Solution 6}
Price risk: $-10 \cdot 1\% = -10\%$ \\
New price predicted by modified duration: NOK 80$ \times 0.90$ = \uuline{NOK 72}



\color{black}
\section*{Problem 7}
A pension plan has an obligation to pay NOK 10,000 once a year during a 10-year period, beginning 5 years from now. The interest rates are flat at 10\%
\begin{enumerate}
     \item[a)] what is the duration of this obligation? 
\end{enumerate}

   

\color{blue}
\subsection*{Solution 7a} 
\noindent The obligation consists of ten annual payments of NOK 10,000, occurring in years 5 through 14. The given annual interest rate is $r=0.10$.

\par\noindent The present value of each payment is:
\[
PV_t=\frac{CF_t}{(1+r)^t},
\]
where $CF_t$ is the payment in year $t$, $r$ is the annual interest rate, and $PV_t$ is the value of that payment today.

\par\noindent The total present value of the obligation is:
\[
P_0=\sum_{t=5}^{14}PV_t
=\sum_{t=5}^{14}\frac{10{,}000}{(1+0.10)^t}
\approx\text{NOK }41{,}968.22.
\]

\par\noindent Macaulay's duration formula is:
\[
D=\sum_{t=5}^{14}t\left[\frac{CF_t(1+r)^{-t}}{P_0}\right],
\]
where $D$ is duration in years and $P_0$ is the total present value calculated above.

\par\noindent For the table, calculate each payment's weight using:
\[
w_t=\frac{PV_t}{P_0}.
\]
The weight $w_t$ is the proportion of the obligation's total present value attributable to the payment in year $t$. Multiply each weight by its payment time $t$, then add the final column to obtain duration. Calculations use unrounded values; the table displays rounded values.

\begin{table}[H]
    \centering
    {\color{blue}
    \begin{tabular}{rrrrrr}
    \toprule
    $t$ & $CF_t$ & $(1+r)^t$ & $PV_t$ & $w_t$ & $w_t\times t$\\
    \midrule
    5  & 10000 & 1.610510 & 6209.21 & 0.147950 & 0.739752\\
    6  & 10000 & 1.771561 & 5644.74 & 0.134500 & 0.807002\\
    7  & 10000 & 1.948717 & 5131.58 & 0.122273 & 0.855911\\
    8  & 10000 & 2.143589 & 4665.07 & 0.111157 & 0.889258\\
    9  & 10000 & 2.357948 & 4240.98 & 0.101052 & 0.909469\\
    10 & 10000 & 2.593742 & 3855.43 & 0.091866 & 0.918655\\
    11 & 10000 & 2.853117 & 3504.94 & 0.083514 & 0.918655\\
    12 & 10000 & 3.138428 & 3186.31 & 0.075922 & 0.911063\\
    13 & 10000 & 3.452271 & 2896.64 & 0.069020 & 0.897259\\
    14 & 10000 & 3.797498 & 2633.31 & 0.062745 & 0.878436\\
    \midrule
    Total & & & 41,968.22 & 1.000000 & \textbf{8.725461}\\
    \bottomrule
    \end{tabular}}
\end{table}

\noindent Therefore, the duration of the obligation is:
\[
D=\sum_{t=5}^{14}t\,w_t
\approx8.725461
\approx\text{\uuline{8.73 years}}.
\]






\noindent
\\
\color{black}
\begin{enumerate}
     \item[b)] If the pension fund intends to use a 5-yr and 20-yr zero coupon bonds to immunise this position, how much should be placed in each bond? 
\end{enumerate}


\color{blue}
\subsection*{Solution 7b} 
\noindent To immunise the obligation, the bond portfolio must have the same present value and duration as the obligation. From part (a), keeping additional decimal places:
\[
P_0\approx\text{NOK }41{,}968.220106,
\qquad
D\approx8.725460512\text{ years},
\]
where $P_0$ is the present value of the obligation and $D$ is its duration.

\par\noindent For a zero-coupon bond, duration equals its remaining maturity:
\[
D_{\mathrm{zero}}=T.
\]
Therefore, the durations of the five-year and twenty-year bonds are:
\[
D_5=5\text{ years},
\qquad
D_{20}=20\text{ years}.
\]

\par\noindent Let $w$ be the proportion of the total investment placed in the five-year bond. The proportion placed in the twenty-year bond is $1-w$. These are market-value weights, not face-value weights.

\par\noindent Portfolio duration is:
\[
D_P=wD_5+(1-w)D_{20}.
\]
To match the obligation's duration, set $D_P=D$:
\[
8.725460512=5w+20(1-w).
\]
Solving for $w$:
\[
w=\frac{20-8.725460512}{20-5}\approx0.751635966,
\]
\[
1-w\approx0.248364034.
\]

\par\noindent\textbf{Market value invested in each bond:}
The amount invested in each bond equals its portfolio weight multiplied by the total amount required:
\[
MV_5=wP_0,
\qquad
MV_{20}=(1-w)P_0,
\]
where $MV_5$ and $MV_{20}$ are the amounts invested today in the five-year and twenty-year bonds.
Therefore:
\[
MV_5=0.751635966\times41{,}968.220106
\approx\text{\uuline{NOK 31,544.82}},
\]
\[
MV_{20}=0.248364034\times41{,}968.220106
\approx\text{\uuline{NOK 10,423.40}}.
\]

\par\noindent\textbf{Face value purchased in each bond:}
For a zero-coupon bond, the relationship between the market value invested and the face value purchased is:
\[
MV=\frac{FV}{(1+r)^T}.
\]
Rearranging:
\[
FV=MV(1+r)^T,
\]
where $FV$ is the face value received at maturity, $MV$ is the amount invested today, $r=0.10$ is the given annual interest rate, and $T$ is the bond's remaining maturity in years.

\par\noindent Using the unrounded market values calculated above:
\[
FV_5=MV_5(1+0.10)^5
\approx\text{\uuline{NOK 50,803.25}},
\]
\[
FV_{20}=MV_{20}(1+0.10)^{20}
\approx\text{\uuline{NOK 70,123.40}}.
\]




\color{black}
\section*{Problem 8}
The 6-month Treasury bill spot rate is 4\% and the 1-yr treasury bill spot rate is 5\%. What is the implied 6-month forward rate for six months from now?\\
\textit{Assume both spot rates are effective annual rates, and express the forward rate on an effective annual basis.}

\color{blue}
\subsection*{Solution 8} 
\noindent Use the Forward Rate Model:
\[
(1+r_T)^T=(1+r_t)^t(1+f_{t,T})^{T-t},
\]
where $r_t$ is the $t$-year spot rate, $r_T$ is the $T$-year spot rate, and $f_{t,T}$ is the effective annual forward rate for the period from time $t$ to time $T$. Both $t$ and $T$ are measured in years from today.

\par\noindent Six months is half a year, so:
\[
t=\frac{6}{12}=0.5,
\qquad
T=1,
\qquad
r_t=0.04,
\qquad
r_T=0.05.
\]

\par\noindent Substituting into the formula:
\[
(1+0.05)^1=(1+0.04)^{0.5}(1+f_{0.5,1})^{1-0.5}.
\]
Rearranging:
\[
(1+f_{0.5,1})^{0.5}
=\frac{(1+0.05)^1}{(1+0.04)^{0.5}}.
\]
Therefore:
\[
f_{0.5,1}
=\left[\frac{(1+0.05)^1}{(1+0.04)^{0.5}}\right]^{1/0.5}-1
\approx\text{\uuline{6.01\%}}.
\]
\noindent This is the effective annual forward rate for the six-month period beginning six months from today.



\noindent
\\
\color{black}
\color{black}
\section*{Problem 9}
Below is a list of prices for zero coupon bonds with par values of NOK 100.
% Please add the following required packages to your document preamble:
% \usepackage{booktabs}
% \usepackage{graphicx}
\begin{table}[H]
\centering
\resizebox{0.25\textwidth}{!}{%
\begin{tabular}{@{}ll@{}}
\toprule
Maturity (Yrs) & Price (\%) \\ \midrule
1              & 94.34      \\
2              & 87.35      \\
3              & 81.64      \\ \bottomrule
\end{tabular}%
}
\label{tab:my-table}
\end{table}

\begin{enumerate}
    \item[a)] An 8.5\% coupon bond making annual coupon will mature in three years. What should the yield to maturity on the bond be?
\end{enumerate}

\color{blue}
\subsection*{Solution 9a} 
The price of the 8.5\%  coupon bond  is NOK 104.02\\
The yield to maturity of the 8.5\%  coupon bond  is 6.97\%


\noindent
\\
\color{black}
\begin{enumerate}
    \item[b)] if at the end of the first year, the yield curve flattens out at 8\%. What will be the 1-yr holding period return on the coupon bond?
\end{enumerate}

\color{blue}
\subsection*{Solution 9b} 
$P_0 = 104.02$\\
$P_1 = 100.89$\\
HPR = 5.16\%


\color{black}
\section*{Problem 10}
The current yield curve for default-free zero-coupon bonds is as follows.
% Please add the following required packages to your document preamble:
% \usepackage{booktabs}
% \usepackage{graphicx}
\begin{table}[H]
\centering
\resizebox{0.35\textwidth}{!}{%
\begin{tabular}{@{}ll@{}}
\toprule
Maturity (Yrs) & Yield to Maturity (\%) \\ \midrule
1              & 10.00     \\
2              & 11.00      \\
3              & 12.00      \\ \bottomrule
\end{tabular}%
}
\label{tab:my-table}
\end{table}
Assuming that the pure expectation hypothesis of the term structure is correct and market expectations are accurate, what will be the yield to maturity on 1-year zero coupon bonds next year?

\color{blue}
\subsection*{Solution 10} 
Under the pure expectation hypothesis, the yield to maturity of 1-year zero coupon bonds next year ($s_2$) is equal to the one-year implied forward rate in year 2 ($f_{1,2}$) = 
12\%
\end{document}
